\subsection{Running example: optimizing the generated F11 TBA}
The shared formula
$\Diamond_{\le5}p\land\Box_{<2}\neg p$ is F11 in the benchmark. Paper~1
uses it to illustrate translation; here we do not repeat that construction.
We begin with the actual TBA emitted by the integrated tool after its
upstream relaxation and reset-completion stages. The checked-in DOT record
\texttt{evaluation/out/F11\_opt\_input.dot} has five locations, eleven
transitions, and two clocks, and is the input to the optimizer measured in
\texttt{evaluation/results.tsv}.

In the first optimization round, contradictory action cubes remove the
transitions that made location~2 reachable. The pass named
\texttt{remove\_unreachable} then removes outgoing transitions from that
location, which now has no incoming edge. Dead-reset elimination and
synchronous-clock merging simplify the remaining timing structure. In the
second round, normalization exposes a reset of the merged-away clock as dead,
and the optimizer removes it. The final TBA has four reachable locations,
seven transitions, and one clock. Export retains seven explicit transitions
and adds two conjunctive location invariants. It also emits a separate
symbolic form with the same accepted language.

Figure~\ref{fig:running-pipeline} shows the first-round and exported
automata. The original generated input is preserved as the DOT artifact
above, rather than replaced by a hand-simplified illustration. This example
demonstrates pass interaction after translation; all translation claims
belong to the earlier translation article~\cite{fares_bodeveix_mtl_translation}.

\begin{figure*}[!t]
\centering
\resizebox{\textwidth}{!}{\begin{tikzpicture}[x=1cm,y=1cm,lab/.style={fmlabel}]
\begin{scope}[xshift=0cm]
  \node[fmnote,font=\small\bfseries] at (2.2,5.3) {After one optimization round};
  \node[fmstate] (l0) at (0.0,4.3) {$0$};
  \node[fmstate] (l1) at (0.0,2.3) {$1$};
  \node[fmstate] (l3) at (4.1,2.3) {$3$};
  \node[fmaccept] (l2) at (0.0,0.0) {$2$};
  \draw[fmedge] (-1.0,4.3) -- (l0);
  \draw[fmedge] (l0) -- node[lab,right] {$\neg p\;/\;\mathsf{rst}(x),\mathsf{rst}(y)$} (l1);
  \draw[fmedge] (l1) edge[loop left,looseness=5] node[lab,left] {$x\le5,\;\neg p$} (l1);
  \draw[fmedge] (l1) -- node[lab,above] {$2\le x\le5,\;\neg p$} (l3);
  \draw[fmedge] (l1) -- node[lab,left] {$2\le x\le5,\;p$} (l2);
  \draw[fmedge] (l3) edge[loop above,looseness=5] node[lab,above] {$x\le5,\;\neg p$} (l3);
  \draw[fmedge] (l3) -- node[lab,pos=.55,below right] {$x\le5,\;p$} (l2);
  \draw[fmedge] (l2) edge[loop right] node[lab,right] {$\top$} (l2);
\end{scope}

\begin{scope}[xshift=8.5cm]
  \node[fmnote,font=\small\bfseries] at (2.2,5.3) {Final TBA};
  \node[fmstate] (r0) at (0.0,4.3) {$0$};
  \node[fmstate,minimum size=11mm] (r1) at (0.0,2.3) {$\begin{array}{c}1\\[-1mm]x\le5\end{array}$};
  \node[fmstate,minimum size=11mm] (r3) at (4.1,2.3) {$\begin{array}{c}3\\[-1mm]x\le5\end{array}$};
  \node[fmaccept] (r2) at (0.0,0.0) {$2$};
  \draw[fmedge] (-1.0,4.3) -- (r0);
  \draw[fmedge] (r0) -- node[lab,right] {$\neg p\;/\;\mathsf{rst}(x)$} (r1);
  \draw[fmedge] (r1) edge[loop left,looseness=5] node[lab,left] {$x\le5,\;\neg p$} (r1);
  \draw[fmedge] (r1) -- node[lab,above] {$2\le x\le5,\;\neg p$} (r3);
  \draw[fmedge] (r1) -- node[lab,left] {$2\le x\le5,\;p$} (r2);
  \draw[fmedge] (r3) edge[loop above,looseness=5] node[lab,above] {$x\le5,\;\neg p$} (r3);
  \draw[fmedge] (r3) -- node[lab,pos=.55,below right] {$x\le5,\;p$} (r2);
  \draw[fmedge] (r2) edge[loop right] node[lab,right] {$\top$} (r2);
\end{scope}
\end{tikzpicture}
}
\caption{F11 after one optimization round and after UPPAAL-oriented export.
The first panel shows the removal of a location with no incoming edges and
the synchronous-clock simplification; the second shows the exported
conjunctive invariants. The full five-location, eleven-transition optimizer
input is the generated artifact cited in the text.}
\label{fig:running-pipeline}
\end{figure*}
